%% Chapter 13A — Operations Runbook
\chapter{Operations Runbook}
\label{ch:operations_runbook}

This chapter defines recurring operating procedures. Print this chapter at semester handover and keep a current copy near the Raspberry Pi.

\section{New Coordinator Handover}

\begin{runbookbox}{Handover Package}
The outgoing coordinator must transfer: dashboard admin access, Firebase project owner contact, Raspberry Pi SSH access method, master key custody location, kiosk station laptop, blank card stock, last known-good firmware binary, and the latest SD card backup.
\end{runbookbox}

\begin{longtable}{L{4cm}L{7cm}L{3cm}}
\toprule
\textbf{Item} & \textbf{Acceptance check} & \textbf{Owner} \\
\midrule
\endhead
Dashboard access & New coordinator can log in and see all machines & Outgoing coordinator \\
Kiosk station & New coordinator issues one test card successfully & Outgoing coordinator \\
RPi access & New maintainer can run \texttt{systemctl status mms-bridge} & Infrastructure maintainer \\
Master key custody & Sealed backup and password-manager entry are verified, not copied casually & Faculty in charge \\
Recovery media & Latest RPi image and known-good firmware binary are readable & Technical maintainer \\
Documentation & This handbook revision is reviewed and dated & New coordinator \\
\bottomrule
\caption{Coordinator handover acceptance checklist}
\label{tab:handover}
\end{longtable}

\section{Semester Transition Checklist}

\begin{enumerate}[leftmargin=1.5em]
    \item Export current users and sessions from Firestore for archival records.
    \item Disable dashboard access for coordinators who are leaving the role.
    \item Add incoming coordinators as Firebase Auth users and Firestore admins.
    \item Count blank cards; reorder before orientation week if stock is below 50.
    \item Run one full access test on every production machine.
    \item Review revoked cards and lost-card incidents from the previous semester.
    \item Confirm the Raspberry Pi SD card backup is less than 45 days old.
    \item Confirm the master key custody record still names current responsible people.
    \item Update the approval page and revision history if the handbook changed.
\end{enumerate}

\section{Daily, Weekly, Monthly Work}

\begin{table}[H]
\centering
\caption{Operations cadence}
\begin{tabular}{L{2.5cm}L{8cm}L{3.5cm}}
\toprule
\textbf{Cadence} & \textbf{Work} & \textbf{Escalate when} \\
\midrule
Daily & Open dashboard; verify no machine is unexpectedly offline and no session is stuck active. & Any relay state disagrees with the physical machine. \\
Weekly & SSH to RPi; check disk, RAM, bridge logs, Mosquitto logs, and card stock. & Repeated bridge exceptions or disk usage above 80\%. \\
Monthly & Make SD backup, review Firebase quota, archive old sessions if needed, run one card issuance drill. & Backup cannot be restored or quota trend is abnormal. \\
Quarterly & Rotate MQTT password if authentication is enabled; review admin list and revoked cards. & Any unknown admin or unexpected credential change appears. \\
\bottomrule
\end{tabular}
\end{table}

\section{Master Key Rotation Runbook}
\label{sec:master_key_rotation}

Rotate the master key only when necessary: suspected leak, coordinator compromise, or planned migration rehearsal. Rotation invalidates every currently issued card.

\begin{enumerate}[leftmargin=1.5em]
    \item Freeze card issuance and announce a maintenance window.
    \item Generate a new 32-byte key on a trusted machine.
    \item Store the new key in the approved password manager and sealed physical backup.
    \item Update \texttt{secrets.h}, \texttt{secrets\_station.h}, and \texttt{secrets.py}.
    \item Re-flash all production nodes and the station writer.
    \item Re-issue cards for all active users.
    \item Keep old revoked-card records for audit, but mark the old key as retired.
    \item Record the rotation reason, date, and approver in the revision log.
\end{enumerate}

\section{Lost Card Procedure}

\begin{enumerate}[leftmargin=1.5em]
    \item Revoke the card immediately in the dashboard.
    \item Confirm revocation reaches all online nodes.
    \item Ask the user where the card was last seen and record the incident.
    \item Issue a replacement card with a new UID.
    \item If the lost card belonged to a coordinator or super admin, also review dashboard access.
\end{enumerate}

\section{Compromised Coordinator Procedure}

\begin{enumerate}[leftmargin=1.5em]
    \item Remove the account from Firebase Auth and Firestore admin collections.
    \item Rotate any shared passwords known to that coordinator.
    \item Review recent dashboard actions: user changes, revocations, OTA triggers, admin additions.
    \item Inspect whether master key custody was exposed. If yes, execute Section~\ref{sec:master_key_rotation}.
    \item Document the incident and the recovery actions.
\end{enumerate}

\section{Disaster Recovery Summary}

\begin{longtable}{L{3.5cm}L{5cm}L{5.5cm}}
\toprule
\textbf{Disaster} & \textbf{Immediate action} & \textbf{Recovery path} \\
\midrule
\endhead
Raspberry Pi failure & Leave nodes powered; sessions buffer locally & Restore latest SD card image or rebuild from deployment guide. \\
Firebase outage & Do not repeatedly restart nodes & Wait for service recovery; restart bridge after Firebase is healthy. \\
ESP32 node failure & Bypass only under electrical supervision & Replace module, flash production firmware, rerun validation. \\
Bad OTA & Stop further rollout & Allow rollback; manually flash failed nodes with known-good binary. \\
Lost master key & Stop issuance & Generate new key, re-flash all nodes, re-issue all cards. \\
Database corruption & Freeze manual edits & Restore machine seed data; reconstruct users from cards if needed. \\
\bottomrule
\caption{Disaster recovery decision table}
\label{tab:dr_summary}
\end{longtable}
